\documentclass{article}
\usepackage{amsmath}
\usepackage{amssymb}
\usepackage{enumitem}

\begin{document}
	
	\section*{Application: Systematic Sampling Using Set Notation}
	
	Consider a population consisting of
	\[
	N = 45,000
	\]
	individuals. A researcher wishes to select a sample of
	\[
	n = 680
	\]
	individuals using \textbf{systematic random sampling}.
	
	The sampling interval is determined by
	\[
	k=\left\lfloor\frac{N}{n}\right\rfloor
	=\left\lfloor\frac{45,000}{680}\right\rfloor
	=66.
	\]
	
	A random starting point is selected from the first sampling interval. Suppose the researcher randomly chooses
	\[
	r=12.
	\]
	
	The first individual selected is therefore the 12th member of the population. Every 66th individual thereafter is included in the sample. The selected individuals are
	\[
	12,\;78,\;144,\;210,\;276,\;\ldots
	\]
	
	This sequence forms an \textbf{arithmetic progression} with
	\begin{itemize}
		\item First term:
		\[
		a_1=12,
		\]
		\item Common difference:
		\[
		d=66.
		\]
	\end{itemize}
	
	Hence, the \(i\)-th sampled individual is given by
	\[
	n_i=12+66(i-1),\qquad i=1,2,\ldots,680.
	\]
	
	The complete sample can therefore be expressed using set notation as
	\[
	S=\left\{12+66(i-1)\mid i=1,2,\ldots,680\right\}.
	\]
	
	Alternatively, using universal quantification,
	\[
	\forall i\in\{1,2,\ldots,680\},\qquad
	n_i=12+66(i-1).
	\]
	
	The final sampled individual is
	\[
	n_{680}
	=12+66(679)
	=44,826,
	\]
	which lies within the population of 45,000 individuals. Thus, the systematic sample consists of exactly 680 individuals selected at equal intervals throughout the population.
	
\end{document}